\begin{lemma}[Bipartite regular-factor cut criterion]
Let $G$ be a finite graph and let $U,W$ be disjoint sets of the same size $n$.
For a natural number $k$, the bipartite graph $G[U,W]$ contains a spanning
$k$-regular subgraph if and only if, for every $A\subseteq U$ and
$B\subseteq W$,
\[
 e_G(A,B)\geq k\max\{0,|A|+|B|-n\}.
\]
Here a spanning $k$-regular subgraph means a subgraph of $G$ in which every
vertex of $U$ has exactly $k$ neighbours in $W$ and every vertex of $W$ has
exactly $k$ neighbours in $U$.
\end{lemma}
\begin{proof}
Suppose first that $F$ is such a subgraph.  For $A\subseteq U$ its vertices
send $k|A|$ edges into $W$.  At most $k|W\setminus B|=k(n-|B|)$ of these
can end outside $B$, because every vertex of $W\setminus B$ has degree $k$.
Consequently
$e_G(A,B)\geq e_F(A,B)\geq k(|A|+|B|-n)$; when the expression in
parentheses is negative the asserted inequality is automatic.

Conversely, form a directed network with a source $s$, a sink $t$, arcs
$s\to u$ of capacity $k$ for $u\in U$, arcs $u\to w$ of capacity $1$ for
each edge $uw$ of $G[U,W]$, and arcs $w\to t$ of capacity $k$ for
$w\in W$.  We prove directly that every $s$--$t$ cut has capacity at least
$kn$.  Write the source side of a cut as
$\{s\}\cup (U\setminus A)\cup B$, where $A\subseteq U$ and $B\subseteq W$.
The arcs crossing it have total capacity
\[
 k|A|+e_G(U\setminus A,W\setminus B)+k|B|.
\]
Apply the assumed inequality to $U\setminus A$ and $W\setminus B$.
Their cardinalities sum to $2n-|A|-|B|$, so the middle term is at least
$k\max\{0,n-|A|-|B|\}$.  If $|A|+|B|\leq n$ the displayed cut capacity is
therefore at least $k(|A|+|B|)+k(n-|A|-|B|)=kn$; if
$|A|+|B|\geq n$, its first and last terms already total at least $kn$.

For completeness, the integral max-flow assertion needed here follows by
the augmenting-path algorithm.  Start with the zero integral flow.  As long
as the residual graph has an $s$--$t$ path, augment by the least residual
capacity on that path.  All capacities and all residual capacities are
integers, so every augmentation is by a positive integer and preserves
integrality.  The flow value strictly increases and is bounded by the total
capacity $kn$ out of $s$, hence the procedure terminates.  If it terminates
below $kn$, let $S$ be the vertices reachable from $s$ in the residual
graph.  Every original arc from $S$ to its complement is saturated and
every original arc in the reverse direction carries zero flow.  Summing
flow conservation over $S\setminus\{s\}$ then says that the flow value is
exactly the capacity of this cut, contradicting the cut bound.  Thus an
integral flow of value $kn$ exists.

All $kn$ units leaving $s$ and entering $t$ are then used.  Since each
$U$--$W$ arc has capacity one and the flow is integral, the arcs carrying
one unit form a subgraph of $G[U,W]$.  Flow conservation and saturation of
the source and sink arcs give degree exactly $k$ at every vertex on both
sides, as required.
\end{proof}
